\section{A local exact solver on cycles with unequal leaf attachments}
\label{sec:op3-local-cycle}

\paragraph{Status and scope.}
This is a \statusproved{} proof draft with an exact implementation and
falsification audit, awaiting independent review. The promised graph is a
simple cycle of unknown length $m\geq3$, with any number $k_i\geq0$ of
degree-one leaves attached to cycle vertex $i$. The $k_i$ may all differ.
The seed is arbitrary. Input consists only of the seed and degree/adjacency
queries; there is no supplied cycle ordering, length, decomposition,
equitable partition or future attachment data. Class verification is not
part of the promise algorithm. General unicyclic graphs with nontrivial
attached trees, and arbitrary cyclic graphs, remain outside this theorem.

\paragraph{Context.}
Root-conditioned scalar homotopy and a symmetry-based pure-cycle solver
already appear in the companion \texttt{delayed\_reflection\_ladder} note.
Here we give the two-tip implementation for unequal leaf attachments,
including its closing-block exception. The proof below is self-contained.

Use the exact-real word model and the degree obstacle problem
\eqref{eq:op3-tree-obstacle}, now with the actual graph degrees. Here
$\lambda>0$ is free; the ACL corollary will set $\lambda=\epsappr$.
All degree replies, adjacency entries, retained records, event searches,
arithmetic and output writes are charged. We first take a cycle seed and
extend to a leaf seed at the end.

\subsection{The continuation parameter and the two exposed ends}

Fix the seed potential $u_v=t\geq0$ and minimize over all other
coordinates. Denote the derivative of this reduced objective by $g(t)$.
It is continuous, strictly increasing and piecewise affine. On each
active face its slope is a positive Schur diagonal. The conditional
coordinates are nondecreasing in $t$, by inverse positivity and the
obstacle comparison principle. The optimum is zero if
$1-\lambda d_v\leq0$; otherwise it occurs at the unique positive root of
$g$. Start at $t=0$, with only the seed exposed.

The two cycle directions from the seed are called arms. Before they meet,
each has an admitted consecutive prefix and one unadmitted candidate. An
unadmitted cycle vertex becomes positive when the sum of its incident
active potentials times $\gamma$ reaches $\lambda d_i$. No other cycle
vertex can activate before one of these candidates. The adjacency row of
a candidate is read only after this event is known to precede the root of
$g$. Its two nonleaf neighbors identify the continuation ports using
charged degree queries. Its other neighbors form a leaf group.

A negative-load leaf attached to a center of potential $x$ has value
\begin{equation}
 u_{\rm leaf}=(\gamma x-\lambda)_+.
 \label{eq:op3-cycle-leaf}
\end{equation}
All $k$ leaves at one center activate at $x=\lambda/\gamma$. Once active,
eliminating them changes its diagonal and load by
\begin{equation}
 d\longleftarrow d-k\gamma^2,\qquad
 b\longleftarrow b-k\gamma\lambda.
 \label{eq:op3-cycle-leaf-fold}
\end{equation}
Reading and storing their labels is charged to the center's degree even
when some leaves are zero in the final answer.

\begin{lemma}[Settled interior prefixes]
\label{lem:op3-cycle-settled}
Before cycle closure, a cycle vertex leaves the active end of an arm only
after its leaf group has activated, or when that group is empty. Therefore
every committed interior record has a fixed quadratic response for the
rest of the continuation.
\end{lemma}
\begin{proof}
A next candidate reached from just this arm requires
$\gamma u_{\rm tip}=\lambda d_{\rm next}\geq2\lambda$.
The tip's leaves activate already at $u_{\rm tip}=\lambda/\gamma$.
Monotonicity prevents later deactivation. The only exception to the
one-sided candidate rule is the last inactive cycle vertex, which is
handled separately below.
\end{proof}

\subsection{Constant-size end records and exact event location}

Eliminate a settled arm prefix outward from the seed. Its current tip has
diagonal $\delta_i>0$ and affine load $\eta_i(t)=a_i t+b_i$, so its
potential is $(a_i t+b_i)/\delta_i$. Initially
$(\delta_i,a_i,b_i)=(d_i,\gamma,-\lambda d_i)$. When committing this tip
and admitting the next cycle vertex $j$, its new record is
\begin{equation}
 \delta_j=d_j-\frac{\gamma^2}{\delta_i},\qquad
 a_j=\frac{\gamma a_i}{\delta_i},\qquad
 b_j=-\lambda d_j+\frac{\gamma b_i}{\delta_i}.
 \label{eq:op3-cycle-append}
\end{equation}
Apply \cref{eq:op3-cycle-leaf-fold} to a tip when its leaf group activates.
All Schur pivots stay positive because they eliminate principal submatrices
of the positive definite Stieltjes matrix.

Maintain two accumulated coefficients over committed records,
\[
 P=\sum_i a_i^2/\delta_i,\qquad
 R=\sum_i a_i b_i/\delta_i.
\]
Let $(\delta_v,b_v)$ denote the seed diagonal/load with any activated seed
leaves folded. While the arms are distinct, the current face derivative is
\begin{equation}
 g(t)=\left(\delta_v-P-\sum_{\rm tips}\frac{a_i^2}{\delta_i}\right)t
       -b_v-R-\sum_{\rm tips}\frac{a_i b_i}{\delta_i}.
 \label{eq:op3-cycle-root-derivative}
\end{equation}
This follows by differentiating the eliminated contributions
$-\eta_i(t)^2/(2\delta_i)$. Each committed record is inserted into $P,R$
once; old records are not refreshed when $t$ changes.

At a state, compute the zero $t_\star$ of the current affine expression
for $g$ and the earliest of at most five prospective events: the seed leaf
group, up to two tip leaf groups, and up to two next cycle candidates.
Every event time follows from a single affine equality using the stored
tip responses. If $t_\star$ is no later than the first event, stop.
Otherwise process that event and repeat. Equal-time events are processed
in any fixed order; a zero at the same time causes stopping before row
exposure. Continuity of $g$ makes this tie rule correct, with zero
coordinates omitted from the output.

\subsection{Closing the cycle without folding inactive leaves}

If the candidate labels of the two arms coincide at $j$, the candidate
receives two-sided load. Write each tip response as $p_a t+q_a$; an empty
arm contributes the seed response $t$. Its activation time is now
\begin{equation}
 t_j=\frac{\lambda d_j-\gamma(q_1+q_2)}
                 {\gamma(p_1+p_2)}.
 \label{eq:op3-cycle-close-time}
\end{equation}
Using the two separate one-sided events here would miss the first
activation. Furthermore, $t_j$ can precede activation of a tip's leaves.
Such a tip cannot yet become a permanently folded interior record.

Keep both current tips and the newly admitted $j$ in a Schur block of
dimension at most three. The rest of each arm already satisfies
\cref{lem:op3-cycle-settled}. The block matrix $\bm S$ has the stored tip
diagonals, diagonal $d_j$, and off-diagonal $-\gamma$ between $j$ and
each nonempty arm tip. Its load is $\bm a t+\bm b$: old tip coefficients
are retained and $a_j=\gamma$ times the number of empty arms. The block
response is
\[
 \bm x(t)=\bm S^{-1}\bm a\,t+\bm S^{-1}\bm b.
\]
Replace the two tip sums in \cref{eq:op3-cycle-root-derivative} by
$\bm a^\top\bm S^{-1}\bm a$ and
$\bm a^\top\bm S^{-1}\bm b$. Any remaining leaf events are obtained from
$x_i(t)=\lambda/\gamma$. There are at most three such groups. Each is
folded by a diagonal/load change inside this same block, and both solves
have dimension at most three. All seed leaves were already settled before
closure by the one-sided first admission. Thus no larger dynamic solve or
search over old vertices is hidden in this step.

At the final $t_\star$, evaluate the tip or closing-block responses. Recover
each arm's committed prefix in reverse order using its stored records:
\begin{equation}
 u_i=\frac{a_i t_\star+b_i+\gamma u_{\rm next}}{\delta_i}.
 \label{eq:op3-cycle-reverse}
\end{equation}
Read each exposed center's leaf list once and use
\cref{eq:op3-cycle-leaf}. Materialize positive coordinates only.

\subsection{Correctness, locality and the approximation guarantee}

\begin{theorem}[Local exact obstacle solve on the promised cyclic family]
\label{thm:op3-local-cycle}
On a simple cycle with arbitrary pendant-leaf counts and an arbitrary
seed, the above algorithm computes the unique obstacle optimum with
$\widetilde O(1+\vol(S))$ charged exact-real word work and storage, where
$S$ is its positive support. Its structural arithmetic, record and query
counts are $O(1+\vol(S))$; the tilde allows deterministic label-dictionary
overhead. It exposes adjacency rows only at vertices in $S$, and
$\lambda\vol(S)\leq1$. It requires no unknown-support preprocessing and
no polynomial factor in $1/\alpha$.
\end{theorem}
\begin{proof}
For a cycle seed, the conditional obstacle has monotone support as $t$
increases. A face starts valid at its activation time, with newly admitted
coordinates zero. Its equations remain valid until the first nonactive
cycle or leaf gate reaches zero. The two-arm description enumerates all
possible cycle gates; disconnected inactive vertices have strictly positive
dual slack. The leaf formula enumerates every remaining type of gate.
\Cref{lem:op3-cycle-settled} justifies committing interiors, while the
small closing block retains all possible unsettled exceptions. The Schur
formulas give the exact reduced derivative and event times on this face.
Its positive slope and continuity certify either the exact optimum before
the next event or the next valid face. At any processed event, $g<0$, so
the eventual root lies strictly later. Monotonicity and strict positive
boundary influence imply that every exposed center is positive finally.
The reverse equations and leaf formulas recover the conditional minimizer
at the root of $g$; all obstacle KKT conditions therefore hold.

There is one admission per exposed cycle vertex and at most one leaf event
per such vertex. Only a constant number of event keys is inspected each
time. Each prefix record is written once and read once in reconstruction;
at most three constant-dimensional leaf updates occur after closure.
Scanning a center and classifying its neighbors costs $O(1+d_i)$, including
all neighbor degree replies. Outputting its leaves and reaccessing stored
labels has the same bound. All exposed centers are in $S$, so these costs
are bounded by $O(1+\vol(S))$, even if some inspected leaf or candidate
labels are outside $S$. A deterministic balanced label dictionary adds at
most logarithmic overhead.

For a leaf seed $s$, first return zero if $\lambda\geq1$. Otherwise its
conditional value at its neighbor $v$ of potential $t$ is
$u_s=1-\lambda+\gamma t>0$. Read the seed's one-entry row and $d_v$.
Eliminating $s$ changes the root data to
\[
 \delta_v=d_v-\gamma^2,\qquad
 b_v=\gamma(1-\lambda)-\lambda d_v.
\]
If $b_v\leq0$, only $s$ is positive, with value $1-\lambda$, and the
root row is not read. Otherwise apply the same cycle-root algorithm with
these modified data, excluding $s$ from its negative leaf group; recover
$s$ by its displayed formula. The additional work is constant and all
exposed rows still belong to positive output coordinates.

Finally sum the optimal equations on $S$. For nonempty support the seed
lies in $S$ and
\[
 1-\lambda\vol(S)
 =\bar\alpha\sum_{i\in S}d_i u_i
   +\gamma\sum_{\substack{\{i,j\}\in E\\i\in S,\ j\notin S}}u_i
 \geq0.
\]
This proves the support-volume bound and completes the accounting.
\end{proof}

\begin{corollary}[The OP3 approximation scale on this graph class]
\label{cor:op3-local-cycle-acl}
Let $\lambda=\epsappr\in(0,1)$ and let
$\widehat{\bm\pi}=\bar\alpha\bm D\bm u$ be the output of
\cref{thm:op3-local-cycle}. For the shared lazy PageRank vector $\bm\pi$,
\[
 \bm0\leq\bm\pi-\widehat{\bm\pi}\leq\epsappr\bm d.
\]
The charged work is $\widetilde O(1/\epsappr)$ in the exact-real word
model on this promised graph family.
\end{corollary}
\begin{proof}
Put $\bm r=\bm e_v-\bm M\bm u$. On positive coordinates
$r_i=\lambda d_i$; on zero coordinates the obstacle conditions and
nonpositive off-diagonals give $0\leq r_i\leq\lambda d_i$.
The empty-support case has the same inequality. Since
$\bm M\one=\bar\alpha\bm d$ and $\bm M^{-1}\geq0$,
\[
 \bm0\leq\bm\pi-\bar\alpha\bm D\bm u
 =\bar\alpha\bm D\bm M^{-1}\bm r
 \leq\lambda\bm d.
\]
Apply the preceding work and volume bounds.
\end{proof}

\paragraph{Measured evidence and limits.}
The exact local routine is \texttt{local\_sun\_solver.py}. The audit
enumerates ordered leaf counts in $\{0,1,2\}$ on cycles of length three
through five, cycle and leaf seeds, and nine parameter pairs: 5,265
comparisons with a separate exact obstacle routine pass. Full-graph KKT
checks are audit-only and excluded from solver counters. Larger asymmetric
cycles demonstrate local stopping without scanning the ambient graph.
This is a family theorem, not arbitrary-graph OP3 and not a claim to realize
every intermediate Wei--Yang active-set gate. Its decisive special property
is that an interior attachment settles before its center leaves the frontier.
Arbitrary attached trees can violate that ordering; identifying a paid
replacement for their changing responses is the next question.

\begin{proposition}[A length-two attachment need not settle; proof draft]
\label{prop:op3-cycle-unsettled-attachment}
The automatic-settling claim fails if pendant leaves are replaced by
arbitrary pendant trees. This refutes that extension of the record rule,
not the existence of an efficient solver on the larger graph class.
\end{proposition}
\begin{proof}
Attach a path $c-a-b$ to a cycle center $c$, where $d_a=2$ and $d_b=1$.
Take $\gamma=3/4$. While $b$ is inactive, the branch response is
$u_a=(\gamma u_c-2\lambda)_+/2$. Thus $a$ activates at
$u_c=8\lambda/3$, and $b$ activates only at $u_c=56\lambda/9$.
A next one-sided cycle neighbor of degree three activates at
$u_c=4\lambda$, strictly between those thresholds. At that point
$u_a=\lambda/2$ and $u_b=0$. Eliminating this branch subtracts $9/32$
from the center diagonal before $b$ activates, but $9/23$ after it
activates. A single fixed quadratic record cannot describe both regimes.

This ordering occurs on an actual source continuation. Use the cycle
$0,1,\ldots,11,0$, add edges $\{1,12\},\{12,13\},\{2,14\}$, seed at
$0$, and $\lambda=1/1000$. The exact conditional solve at
$u_0=39/2000$ gives $u_1=1/250$, $u_2=0$, $u_{12}=1/2000$, $u_{13}=0$
and reduced source derivative $-10639/11000<0$. At $u_0=1/20$ it gives
$u_2=307/118500>0$, $u_{13}=89/31600>0$ and source derivative
$-2147709/2325760<0$. The full vectors and exact conditional
solves are recorded in the audit. Both events therefore precede
the final root, and the old branch changes after the next cycle admission.
\end{proof}
